\documentclass[conference,10pt]{IEEEtran}
\IEEEoverridecommandlockouts

\usepackage{cite}
\usepackage{amsmath,amssymb,amsfonts}
\usepackage{algorithmic}
\usepackage{graphicx}
\usepackage{textcomp}
\usepackage{xcolor}
\usepackage{booktabs}
\usepackage{url}
\usepackage{hyperref}
\usepackage{enumitem}
\usepackage{bbm}
\usepackage{microtype}
\usepackage{balance}

\def\BibTeX{{\rm B\kern-.05em{\sc i\kern-.025em b}\kern-.08em
    T\kern-.1667em\lower.7ex\hbox{E}\kern-.125emX}}

\begin{document}

\title{Revised Plan v3: Addressing Reviewer Concerns for the Dual-Mode
Quantum MRF Sampling Manuscript}

\author{
\IEEEauthorblockN{Arul Rhik Mazumder}
\IEEEauthorblockA{\textit{School of Computer Science} \\
\textit{Carnegie Mellon University}\\
Pittsburgh, PA, USA \\
arulm@cs.cmu.edu}
}

\maketitle

\begin{abstract}
This document is a revision plan, not a research paper. It enumerates the
substantive concerns raised by the IEEE QCE peer review of the manuscript
\emph{``A Dual-Mode Quantum Framework for Sampling Discrete Markov Random
Fields''}, classifies each as \textbf{must-fix}, \textbf{should-fix}, or
\textbf{defer-to-future-work}, and specifies the concrete changes to text,
experiments, and citations that are required for resubmission. The plan is
structured so that every reviewer concern maps to (a) one or more specific
edits in the manuscript, (b) any new experiments needed to support the edit,
and (c) an acceptance criterion that determines when the revision is
complete. Twelve revision items are tracked in total. If items
R1--R8 are completed, the manuscript should be competitive at IEEE QCE.
Items R9--R12 are deferred to a follow-up paper.
\end{abstract}

\begin{IEEEkeywords}
Revision plan, peer review response, quantum sampling, MRF, IEEE QCE,
reproducibility.
\end{IEEEkeywords}

\section{Scope and Triage}

The reviewer raised eight major concerns and several minor ones. We classify
each by severity and effort:

\begin{itemize}[leftmargin=*]
    \item \textbf{Must-fix before resubmission (R1--R5).} Issues that, if
    left unaddressed, will cause rejection regardless of other merits.
    These are: incorrect citation of the QCGM paper, the tautological
    framing of the headline ESS result, the misleading framing of the
    variational compression result, the incorrect framing of the null
    clique-entanglement finding, and several internal inconsistencies and
    presentation errors.
    \item \textbf{Should-fix to be competitive (R6--R8).} Issues that
    individually are survivable but collectively determine whether the
    paper reads as a careful contribution or a preliminary report. These
    are: adding at least one stronger classical baseline, normalising the
    ESS comparison by wall-clock cost, and broadening the citation
    landscape to cover Born machines, Wittek--Gogolin, and recent
    QeMCMC follow-ups.
    \item \textbf{Defer to future work (R9--R12).} Items that the reviewer
    would like to see but that are out of scope for the current artifact.
    These are: the full ancilla-based QCGM implementation, a hardware run,
    multi-seed coverage at $n > 12$, and the marginal-matching variational
    objective. Each is acknowledged in the limitations section with a
    concrete handoff to a follow-up paper.
\end{itemize}

The remainder of this document walks each item.

\section{Must-Fix Revisions}

\subsection{R1 --- Correct the QCGM citation}

\textbf{Concern.} Reference~[7] in the submitted manuscript reads
``N.~Piatkowski, `Quantum circuits for discrete graphical models,' in
Proc.~IEEE QCE, 2022.'' This is wrong on three counts: the paper has two
authors (Piatkowski and Zoufal), the venue is \emph{Quantum Machine
Intelligence} (Springer), and the publication year is 2024 (the 2022
date refers to the arXiv preprint, arXiv:2206.00398). The author's own
GitHub README cites the paper correctly, so this is a paper-internal
inconsistency.

\textbf{Action.} Replace the bibliography entry with:
\begin{quote}\small
N.~Piatkowski and C.~Zoufal, ``On quantum circuits for discrete graphical
models,'' \emph{Quantum Machine Intelligence}, vol.~6, art.~57, 2024.
[Online]. Available: \url{https://doi.org/10.1007/s42484-024-00175-y}.
Originally available as arXiv:2206.00398, 2022.
\end{quote}
Update every in-text reference to QCGM accordingly. The text in
\S\,Related Work that paraphrases the original method should be checked
to ensure it accurately reflects the published version (which differs
from early preprints in some details of the success-probability bound).

\textbf{Acceptance criterion.} A literature-search reviewer can locate the
paper from the citation as printed.

\subsection{R2 --- Reframe the ESS result honestly}

\textbf{Concern.} The headline ``$17.38\times$ ESS advantage'' compares
i.i.d. quantum samples against single-site random-scan Gibbs. This is
structurally tautological: amplitude-encoded samples have $\tau \approx 1$
\emph{by construction}, and single-site Gibbs on loopy graphs is the
weakest reasonable MCMC baseline. The reviewer correctly notes that the
result, as currently framed, restates ``i.i.d. samples are uncorrelated.''

\textbf{Action.} Three coordinated edits:
\begin{enumerate}[leftmargin=*]
    \item \textbf{Abstract and introduction.} Replace ``quantum sample-quality
    advantage'' with ``ESS gap relative to single-site Gibbs.'' Make explicit
    that the comparison is informative about the cost of mixing in vanilla
    MCMC, not about a quantum-vs-classical advantage in absolute terms.
    \item \textbf{New baseline (see also R6).} Add a block-Gibbs comparator
    that updates an entire clique at a time; report ESS ratios for both
    baselines. The expectation is that block Gibbs will substantially
    narrow the gap on the chain and two-clique families and reduce or
    eliminate it on the barbell families. Whatever the outcome, report it.
    \item \textbf{Limitations subsection.} Promote the existing
    ``Comparison with advanced MCMC'' bullet to a full subsection
    (\S\,Discussion) titled \emph{``What the ESS comparison does and does
    not show''}, stating that the experiment quantifies the cost of using
    the simplest MCMC variant and does not rule out classical samplers
    with comparable autocorrelation.
\end{enumerate}

\textbf{Acceptance criterion.} The abstract no longer claims ``quantum
sample-quality advantage'' as a primary contribution. The discussion
makes it possible for a reader to estimate, within an order of
magnitude, how much of the $17\times$ gap survives a stronger baseline.

\subsection{R3 --- Recast the variational compression as a negative result}

\textbf{Concern.} Table~III shows MPS dominating VQC at every measured
size, with the absolute fidelity gap widening from $n=8$ ($-0.68$) to
$n=10$ ($-0.75$) before slightly narrowing at $n=12$ ($-0.71$). Calling
this ``compression of $113.8\times$'' is misleading: compression at
$F=0.165$ is not compression in any operationally useful sense.

\textbf{Action.}
\begin{enumerate}[leftmargin=*]
    \item Rename the contribution in the introduction from ``a variational
    compression scheme'' to ``a characterization of the parameter--fidelity
    trade-off for hardware-efficient ans\"atze on MRF distributions.''
    \item Drop the unqualified compression-ratio numbers from the abstract.
    Where they appear in the body, always pair them with the corresponding
    fidelity value, e.g. ``$113.8\times$ parameter reduction at
    $F_{\mathrm{VQC}} = 0.165$.''
    \item In the discussion, state explicitly: ``At every size measured,
    a depth-3 hardware-efficient ansatz with $nd$ parameters trained for
    $\sim 75$ steps fails to match an MPS with bond dimension matched to
    the same parameter budget. Whether this gap closes at larger $n$ or
    deeper $d$ is an open empirical question; the present data do not
    support a positive crossover.''
    \item Add a paragraph in \S\,Future Directions identifying the
    candidate fixes (deeper circuits with sensible initialization;
    natural-gradient or quantum-natural-gradient optimization;
    marginal-matching objective; clique-aware initial-state preparation),
    and a clear statement that establishing a crossover is the central
    open question.
\end{enumerate}

\textbf{Acceptance criterion.} A reader who reads only the abstract and
the table captions cannot mistake the compression result for a positive
finding.

\subsection{R4 --- Demote the clique-entanglement contribution}

\textbf{Concern.} The clique-vs-linear effect is $\Delta F = -2.28\times
10^{-4}$ with $R^2 = 6.55\times 10^{-4}$. This is a null result, not a
characterization. Listing it as a numbered novel contribution overstates
what the data support.

\textbf{Action.}
\begin{enumerate}[leftmargin=*]
    \item Remove ``problem-aware (clique-based) entanglement'' from the
    numbered contributions list in the introduction.
    \item Move the entanglement sweep to \S\,Results as a diagnostic, with
    a one-sentence summary: ``In our 36-instance sweep, clique
    entanglement provided no measurable advantage over linear entanglement
    at the sizes tested ($n \le 14$).''
    \item In Fig.~3 caption, drop the slope value ($2.58$) entirely. With
    $R^2 = 0.003$ the slope is not interpretable; reporting it invites
    misreading.
    \item Add to limitations: ``Larger graphs and denser entanglement
    structures than those tested here would be needed to detect any
    benefit of clique-aligned entanglement.''
\end{enumerate}

\textbf{Acceptance criterion.} The introduction's contribution list
contains only contributions that the experimental data substantiate.

\subsection{R5 --- Fix internal inconsistencies and presentation errors}

\textbf{Concern.} The reviewer flagged several issues that individually
are minor but collectively erode credibility.

\textbf{Action.} Apply the following edits:
\begin{enumerate}[leftmargin=*]
    \item \textbf{Parameter-counting convention.} Tables~II and~III use
    different conventions ($2nd$ vs $nd$). Pick one (recommendation: $nd$,
    since the matched-budget crossover table is the more important
    artifact), recompute Table~II rows under that convention, and add a
    single explanatory sentence in \S\,Methods clarifying that $nd$
    counts only $R_Y$ angles in the alternating layer scheme.
    \item \textbf{Grid graphs.} Section~I-B mentions $n=9$ and $n=16$ grid
    graphs ``representative of image-segmentation workloads,'' but they
    never appear in the experiments. Either add at least one row of grid
    results to Table~III, or remove the claim from \S\,I-B and replace it
    with the actual experimental coverage (chain, barbell, barbell-path,
    Erd\H{o}s--R\'enyi, two-clique).
    \item \textbf{Reference~[2] (S\"oding 2005).} The cited paper is on
    HMM-HMM protein homology detection, not MRFs. Replace with Kamisetty
    et al., \emph{PNAS} 2013 (``Assessing the utility of coevolution-based
    residue-residue contact predictions in a sequence- and
    structure-rich era'') for the computational-biology motivation.
    \item \textbf{``Exact samples.''} In the abstract and \S\,I-B,
    reserve the phrase ``exact samples'' for the amplitude-encoding path
    only. The variational path produces samples from $\hat P_{\theta}$,
    not $P_{\theta}$.
    \item \textbf{Table~II inversion.} Note in the table caption that
    fidelity at $n=14$ ($0.438$) is below the $n=6$ enumeration-competitive
    point ($0.519$), and explain in one sentence what this means: with
    fixed depth $d=3$, the ansatz's expressive power does not keep pace
    with the growing distribution dimension.
    \item \textbf{Reproducibility timestamp.} The current note records
    artifact regeneration on 2026-04-25, the day of submission. Add a
    sentence stating that the seeds used for each experiment are listed
    in the repository's \texttt{experiments/seeds.json}, and include the
    seed list as supplementary material.
\end{enumerate}

\textbf{Acceptance criterion.} A second reading pass produces no
conventions clashes, no orphaned claims (grid graphs), and no obviously
unsupported statistical statements.

\section{Should-Fix to Be Competitive}

\subsection{R6 --- Add at least one stronger classical baseline}

\textbf{Concern.} Without a stronger MCMC baseline, the headline ESS
result is uninformative.

\textbf{Action.} Implement and run \emph{block Gibbs} as the new primary
classical baseline. Block Gibbs updates each maximal clique jointly via
exact enumeration over the clique's $2^{|C|}$ assignments conditional on
its boundary. Implementation effort: roughly one day, since the diagonal
of $H_\theta$ is already computed for amplitude encoding and the same
machinery enumerates clique-conditional probabilities. Run the same 60
instances with the block-Gibbs comparator alongside single-site Gibbs.
Report the new ESS ratios as an additional column in the ESS table.

If time permits, add a Swendsen--Wang comparator for the two-clique and
barbell families. Swendsen--Wang is the standard ``hard'' baseline for
clustered Ising-like distributions and is the comparison that an
informed reviewer will ask about.

\textbf{Acceptance criterion.} The ESS table reports Quantum-vs-Single-Site
\emph{and} Quantum-vs-Block ratios for all 60 instances. The text
discusses how much of the original $17\times$ gap survives.

\subsection{R7 --- Wall-clock-normalized comparison}

\textbf{Concern.} The current ESS metric counts samples but ignores the
$O(2^n)$ classical preprocessing required for amplitude encoding. The
relevant operational quantity is \emph{ESS per second of total wall
clock}, not ESS per sample.

\textbf{Action.} Add a new table reporting, for each method
(amplitude-encoded quantum, single-site Gibbs, block Gibbs, and where
applicable Swendsen--Wang):
\begin{itemize}[leftmargin=*]
    \item Total wall-clock time including preprocessing.
    \item Number of effective samples produced.
    \item ESS per second.
\end{itemize}
The expectation is that block Gibbs will dominate amplitude encoding on
ESS-per-second at the sizes where preprocessing is the binding cost
($n \ge 10$). State this honestly; it does not undermine the paper, but
it does refine the claim.

\textbf{Acceptance criterion.} The discussion contains a clear statement
of the regime in which amplitude encoding is competitive on a wall-clock
basis (e.g., ``small $n$ with many independent samples per preprocessing
batch'') and the regime in which it is not.

\subsection{R8 --- Broaden the citation landscape}

\textbf{Concern.} The current bibliography omits work the reviewer
considers prior art:
\begin{itemize}[leftmargin=*]
    \item Liu \& Wang, \emph{Phys.~Rev.~A} 98, 062324 (2018), ``Differentiable
    learning of quantum circuit Born machines.''
    \item Benedetti et al., \emph{npj Quantum Inf.} 5, 45 (2019),
    ``A generative modeling approach for benchmarking and training shallow
    quantum circuits.''
    \item Wittek \& Gogolin, \emph{Sci.~Rep.} 7, 45672 (2017),
    ``Quantum enhanced inference in Markov logic networks.''
    \item Ferguson \& Wallden, \emph{Phys.~Rev.~Research} 7, 013231 (2025),
    ``Quantum-enhanced MCMC for systems larger than your quantum computer.''
\end{itemize}

\textbf{Action.} Add the four citations. Wittek--Gogolin in particular
should be discussed in \S\,Related Work, because it is closer to MRF
sampling than the Layden et al.\ Ising-MCMC work that currently anchors
that paragraph. Liu--Wang and Benedetti et al.\ should be cited at the
point where the variational-circuit-as-Born-machine framing is introduced
in \S\,III-B.

\textbf{Acceptance criterion.} The four references appear and are each
discussed in at least one sentence of the body text.

\section{Defer to Future Work}

\subsection{R9 --- Full ancilla-based QCGM implementation}

The reviewer notes, correctly, that the simplified amplitude-encoding
path discards the very property (exponentially-better-than-classical
preprocessing in the favorable parameter regime) that motivates the
QCGM framework. The author estimated $2$--$4$ weeks of engineering. This
is the central technical extension and is the natural follow-up paper.

\textbf{Action for current submission.} Add a paragraph in \S\,Future
Directions that scopes this work explicitly: ``A follow-up effort,
estimated at four weeks, will implement the ancilla-based repeat-until-
success construction of Piatkowski and Zoufal in our codebase, replacing
the $O(2^n)$ classical preprocessing with the $O(\mathrm{poly}(|\mathcal{C}|))$
quantum sampling primitive. That work will determine whether the
amplitude-encoding fidelity guarantees demonstrated here transfer to the
full quantum-advantage regime.''

\subsection{R10 --- Hardware run}

The reviewer notes that for a venue called \emph{Quantum Computing and
Engineering}, all-simulator results are a hard sell.

\textbf{Action for current submission.} Add to limitations and to
future directions an explicit hardware-deployment plan: target IBM
Heron (or comparable) for a pilot run at $n \in \{4, 6, 8\}$ on chain
graphs, with readout error mitigation and zero-noise extrapolation.
Acknowledge that no hardware data are reported in the present paper.
If feasible before the final submission deadline, include even a single
$n=4$ chain hardware result; the absolute fidelity does not need to be
high to demonstrate that the pipeline runs end-to-end on a real device.

\subsection{R11 --- Multi-seed coverage at $n > 12$}

The MPS scaling table reports single-seed runs at $n \in \{12, 24, 32, 40\}$.
The reviewer correctly notes that single-seed results at the edge of the
study cannot be used to argue scaling trends.

\textbf{Action for current submission.} Either run three additional seeds
per $(n, \chi)$ cell in Table~VII (estimated cost: a few hours of
BlueQubit time), or, if cost is prohibitive, add explicit error bars
based on the within-instance variance of the single seed and clearly
label the rows as preliminary. The former is preferred.

\subsection{R12 --- Marginal-matching objective}

The current variational training uses KL divergence end-to-end. A
clique-marginal-matching objective could exploit the MRF factorization
explicitly and is likely to improve fidelity by a meaningful amount.

\textbf{Action for current submission.} List as a future direction with
an estimated effort of one week and a concrete prediction (the
$0.1$--$0.15$ fidelity gain mentioned in the original future-directions
list). No experimental work in the current revision.

\section{Cross-Cutting Edits}

\subsection{Restructure the contributions list}

The introduction's numbered contribution list, after R2 and R4, should
read:
\begin{enumerate}[leftmargin=*]
    \item A reproducible benchmark of amplitude encoding for discrete
    MRFs at $n = 8, 10, 12$, validating bit-ordering, normalization, and
    the diagonal-only memory optimization to near-unit fidelity on
    BlueQubit SV.
    \item A graph-family-stratified characterization of the ESS gap
    between amplitude-encoded i.i.d. sampling and two MCMC baselines
    (single-site Gibbs and block Gibbs), reported on a wall-clock-normalized
    basis across 60 instances.
    \item An empirical study of hardware-efficient ans\"atze as
    distribution compressors at $n \in \{8, 10, 12\}$, finding that depth-3
    circuits with $nd$ parameters fail to match MPS at matched parameter
    budgets, and identifying the candidate fixes whose evaluation is the
    natural follow-up.
    \item An open-source codebase with executable artifact-verification
    scripts.
\end{enumerate}

This list is shorter than the submitted version, but each item is
defensible from the data.

\subsection{Title}

The current title---``A Dual-Mode Quantum Framework for Sampling
Discrete Markov Random Fields''---promises more than the paper now
delivers. Two candidates:
\begin{itemize}[leftmargin=*]
    \item \emph{``A Reproducible Benchmark of Quantum Amplitude Encoding for
    Discrete Markov Random Field Sampling''} --- conservative, accurate.
    \item \emph{``Quantum Amplitude Encoding for Discrete MRFs: An ESS
    Characterization Against Gibbs Baselines''} --- centers the headline
    finding.
\end{itemize}
Either is preferable to the current title once the variational result
is recast as in R3.

\subsection{Status table}

The status table (Table~IX) is one of the strongest features of the
manuscript and should be preserved. Update it to reflect the new state
after R2--R4 and R6--R7:
\begin{itemize}[leftmargin=*]
    \item ESS characterization: now ``Complete with two baselines.''
    \item VQC compression: now ``Negative result at $n \le 12$; crossover
    open.''
    \item Clique entanglement: now ``Null at tested scales; demoted from
    contribution.''
    \item Hardware run: ``Deferred to follow-up.''
\end{itemize}

\section{Estimated Effort and Timeline}

The must-fix items (R1--R5) are pure-text edits and require roughly
two days of focused writing, with R5 requiring a careful pass to
unify conventions across tables. The should-fix items (R6--R8)
require new experiments:
\begin{itemize}[leftmargin=*]
    \item R6 (block Gibbs implementation and 60-instance run): one day.
    \item R7 (wall-clock-normalized table): half a day, since the timing
    instrumentation is already in place.
    \item R8 (citation additions): half a day.
\end{itemize}
Total revision effort: approximately one work week. R9--R12 are
out of scope and deferred.

\section{Acceptance-Probability Estimate (Self-Assessed)}

The reviewer estimated the current paper at 10--15\% acceptance, and
estimated that completing the must-fix and should-fix items would put
the revised paper at 45--55\%. We accept that estimate. The dominant
remaining risk after R1--R8 is the absence of hardware results (R10),
which a strict QCE reviewer may still consider disqualifying. If the
hardware pilot in R10 can be executed before the final deadline, even
at $n = 4$, the paper moves into the upper end of that range.

\section{Summary}

This revision plan converts the eight major reviewer concerns into
twelve concrete revision items, of which five are pure-text edits
(R1--R5), three require additional experiments at modest cost (R6--R8),
and four are deferred to a follow-up effort (R9--R12). The plan
preserves the strengths the reviewer identified---honest reporting,
reproducible code, the status table---and removes the framings that
the reviewer correctly flagged as overclaims. If executed, the revised
manuscript will be a defensible benchmark paper rather than an
overclaimed result paper, which is the appropriate genre for the
current state of the work.

\end{document}